\documentclass[12pt]{evaluasep}
\grado{1$^\circ$ de Secundaria}
\cicloescolar{2025-2026}
\materia{Olimpiada de Conocimientos 2026}
\unidad{1}
\title{Propuesta para Examen}
\author{Profesores de Secundaria}
\begin{document}%
% \begin{multicols}{2}
	\tableofcontents
% \end{multicols}
\vfill
\newpage
\begin{questions}
		\section{Lenjuages}
                \subsection{Lengua Materna}

				\question Indica cuáles son las fases estructurales de las mesas de diálogo:
    \begin{choices}
        \choice Principio, conversación y conclusión.
        \CorrectChoice Apertura, desarrollo y cierre.
        \choice Principio, continuación y diálogo.
        \choice Conversación, diálogo y cierre.
    \end{choices}
    

    \question ¿Qué es un Decálogo?
    \begin{choices}
        \choice Una serie conformada por 7 reglas o normas básicas de conducta.
        \choice Una serie estructurada por 20 estatutos o leyes secundarias.
        \CorrectChoice Una serie integrada por 10 reglas o principios básicos.
        \choice Una serie de 5 recomendaciones opcionales dentro de un manual de usuario.
    \end{choices}
    

    \question ¿Cuáles son las funciones principales de los medios de comunicación masiva?
    \begin{choices}
        \choice Educar, vender, informar, promover y persuadir.
        \choice Enseñar, educar, contener, vender y promover.
        \choice Informar, educar, vender, entretener y promover.
        \CorrectChoice Informar, educar, entretener, promover y persuadir.
    \end{choices}
    

    \question Selecciona la opción que contenga ejemplos reales de medios de comunicación masiva:
    \begin{choices}
        \CorrectChoice Periódicos, televisión, radio e internet.
        \choice Revistas académicas cerradas, televisión local, radioafición y gacetas internas.
        \choice Cómics independientes, periódicos murales de aula, internet y televisión satelital.
        \choice Internet de banda ancha, periódicos históricos, revistas descatalogadas y cómics clásicos.
    \end{choices}
    

    \question ¿Qué elementos es común e indispensable incluir dentro de las cédulas museográficas?
    \begin{choices}
        \CorrectChoice Título de la obra, nombre del artista, año de creación, técnica empleada y texto explicativo.
        \choice Título de la obra, descripción de cómo se hizo, para quién estaba destinada y por qué se realizó.
        \choice Nombre del artista, datos del ayudante de taller, lugar exacto de origen y texto crítico descriptivo.
        \choice Técnica artística utilizada, texto poético libre, nombre del museo, año de inauguración y nombre del artista.
    \end{choices}
	

				    \subsection{Lengua extranjera}

\question Which of the following sentences correctly uses the present perfect tense to express an action that started in the past and continues to the present?
    \begin{choices}
        \choice I am living in this city since 2015.
        \CorrectChoice I have lived in this city since 2015.
        \choice I have been living in this city for 2015.
        \choice I live in this city from 2015.
    \end{choices}
   
\newpage
    \question Identify the sentence with the correct order of adjectives:
    \begin{choices}
        \choice She bought a red beautiful leather Italian bag.
        \CorrectChoice She bought a beautiful red Italian leather bag.
        \choice She bought an Italian beautiful red leather bag.
        \choice She bought a leather beautiful red Italian bag.
    \end{choices}
   

    \question Read the following sentence: \textit{"Despite the torrential downpour, the match continued until the final whistle."} What can be logically inferred from this statement?
    \begin{choices}
        \choice The match was canceled completely due to bad weather.
        \choice The players stopped playing because of the heavy rain.
        \CorrectChoice The match finished completely even though it was raining heavily.
        \choice The referee blew the whistle to pause the game temporarily.
    \end{choices}
   

    \question If someone uses the idiom \textit{"It is raining cats and dogs"}, they mean that:
    \begin{choices}
        \choice Animals are literally falling from the sky.
        \CorrectChoice It is raining very heavily.
        \choice There is a storm combined with many stray animals on the street.
        \choice The weather is highly unpredictable and strange.
    \end{choices}
   

    \question Choose the correct option to complete the second conditional sentence: \textit{"If I \underline{\hspace{2cm}} a superpower, I \underline{\hspace{2cm}} the ability to fly."}
    \begin{choices}
        \choice have / choose
        \CorrectChoice had / would choose
        \choice will have / will choose
        \choice would have / choose
    \end{choices}

\newpage

\section{Saberes y pensamiento científico}
\subsection{Matemáticas}

\question ¿Cuál es la suma de los dígitos del múltiplo común más pequeño de 9, 10 y 11 que es mayor a cero?
    \begin{choices}
        \choice 9
        \CorrectChoice 18
        \choice 27
        \choice 99
    \end{choices}


    \question Un reloj se atrasa 4 minutos cada 12 horas. Si se sincroniza a la perfección un lunes a las $8:00$ a.m., ¿qué hora real será cuando el reloj marque las $8:00$ a.m. del viernes de esa misma semana?
    
	\begin{choices}
        \choice $7:28$ a.m.
        \choice $7:32$ a.m.
        \CorrectChoice $8:32$ a.m.
        \choice $8:28$ a.m.
    \end{choices}
	

    \question Si $x$ e $y$ son números primos distintos tales que $x + y = 31$, ¿cuál es el valor del producto $x \cdot y$?
    \begin{choices}
        \choice 58
        \CorrectChoice 58 %(Dado que el único primo par es 2, $x=2$ e $y=29$, su producto es 58).
        \choice 114
        \choice No existe tal par de números primos.
    \end{choices}
	

    \question En un cajón hay 12 calcetines negros, 8 azules y 6 blancos. Si extraes calcetines a oscuras, ¿cuál es el mínimo número de calcetines que debes sacar para garantizar que tienes al menos un par del mismo color?
    \begin{choices}
        \choice 2
        \choice 3
        \CorrectChoice 4
        \choice 7
    \end{choices}
	

    \question El área de un rectángulo es de 48 cm² y su perímetro es de 28 cm. ¿Cuál es la longitud de la diagonal de dicho rectángulo?
    \begin{choices}
        \choice 8 cm
        \CorrectChoice 10 cm
        \choice 12 cm
        \choice 14 cm
    \end{choices}
	

\subsection{Biología}

\question ¿Cuál es la función principal del ADN en los seres vivos?
  
\begin{choices}
        \choice Proteger al cuerpo contra enfermedades en el torrente sanguíneo.
        \choice Producir energía química directa para que las células funcionen.
        \choice Transportar oxígeno hacia los tejidos periféricos de los animales.
        \CorrectChoice Guardar la información genética y dirigir la producción de proteínas.
    \end{choices}
    

    \newpage

    \question ¿Cuál es una diferencia importante entre las células procariotas y las eucariotas?
  
    \begin{choices}
        \choice Las procariotas producen notablemente más energía metabólica que las eucariotas.
        \CorrectChoice Las procariotas no tienen núcleo definido, mientras que las eucariotas sí tienen núcleo y organelos.
        \choice Las eucariotas carecen de membrana celular, pero las procariotas sí poseen una capa protectora rígida.
        \choice Las procariotas siempre presentan dimensiones mayores en comparación con las eucariotas.
    \end{choices}
    

    \question ¿Qué es una cadena alimentaria?
 
    \begin{choices}
        \choice Una red de carreteras geográficas que conecta diferentes ecosistemas contiguos.
        \choice El conjunto específico de organismos vegetales que producen oxígeno en un hábitat.
        \choice La lista exhaustiva de animales exóticos que habitan en un mismo ecosistema protegido.
        \CorrectChoice La ruta de transferencia de energía y nutrientes a través de los diferentes organismos en un ecosistema.
    \end{choices}
    

    \question ¿Qué tipo de células defienden al cuerpo de las infecciones?
   
    \begin{choices}
        \choice Los glóbulos rojos o eritrocitos.
        \CorrectChoice Los glóbulos blancos o leucocitos.
        \choice Las células del sistema nervioso o neuronas.
        \choice Las células musculares o miocitos.
    \end{choices}
    

    \question ¿Qué es una red trófica?
   
    \begin{choices}
        \choice Una cadena de montañas o fallas geológicas que delimita ecosistemas aislados.
        \choice El conjunto de productores primarios que captan de forma exclusiva la energía solar.
        \CorrectChoice El conjunto de cadenas alimentarias interconectadas que muestran cómo se transfiere la energía y los nutrientes entre los organismos de un ecosistema.
        \choice La lista de animales en peligro de extinción que residen en un mismo territorio o reserva natural.
    \end{choices}
	

\section{De lo Humano a lo comunitario}
\subsection{Tecnología}

\question ¿Qué es Google Classroom y para qué sirve principalmente?
    \begin{choices}
        \choice Es un editor de fotografías digitales avanzado que sirve para diseñar materiales educativos de alta definición.
        \CorrectChoice Es una plataforma educativa de Google que permite a maestros y alumnos compartir tareas, archivos, avisos y actividades de manera organizada y en línea.
        \choice Es una herramienta de videoconferencias masivas que sustituye la necesidad de enviar archivos o actividades por correo.
        \choice Es un disco duro físico que se conecta a las computadoras escolares para almacenar respaldos de bases de datos de control escolar.
    \end{choices}
    

    \question ¿Cómo se crea una carpeta o un archivo nuevo directamente en la plataforma Google Classroom?
    \begin{choices}
        \choice Modificando la configuración de la cuenta personal en el menú de privacidad de Google.
        \choice Descargando un programa externo e ingresando el código secreto provisto por la dirección escolar.
        \CorrectChoice Entrando a una tarea o apartado de ``Trabajo en clase'', seleccionando la opción ``Agregar o crear'' y eligiendo el tipo de archivo deseado.
        \choice Enviando un mensaje de soporte técnico directo a los administradores de la plataforma educativa.
    \end{choices}
    
    \question ¿Qué pasos se deben seguir para subir correctamente un archivo ya existente a Google Classroom?
    \begin{choices}
        \choice Compartir el archivo en las redes sociales institucionales de la escuela y etiquetar al profesor de la materia.
        \CorrectChoice Abrir la tarea correspondiente, hacer clic en ``Agregar archivo'', seleccionar el documento desde el dispositivo o Google Drive y presionar ``Entregar''.
        \choice Enviar un correo electrónico con el documento adjunto a cada uno de los integrantes del grupo escolar.
        \choice Arrastrar el documento directamente al motor de búsqueda general de Google sin abrir la sesión.
    \end{choices}
    

    \question ¿Por qué es de suma importancia organizar de forma correcta los archivos dentro de Google Classroom?
    \begin{choices}
        \choice Porque aumenta de manera automática la velocidad de la conexión a internet de los dispositivos domésticos.
        \choice Porque permite instalar aplicaciones y videojuegos adicionales sin ocupar espacio en la memoria interna.
        \CorrectChoice Porque facilita encontrar tareas, mantener el trabajo ordenado, evitar pérdidas de información y mejorar la comunicación entre alumnos y docentes.
        \choice Porque es un requisito de seguridad indispensable para evitar que la cuenta de correo electrónico institucional sea suspendida.
    \end{choices}
    

    \question ¿Qué son los Documentos de Google y cuál es su función primordial?
    \begin{choices}
        \choice Son plantillas de diseño gráfico destinadas exclusivamente al renderizado de animaciones en tercera dimensión.
        \choice Son bases de datos complejas para el cálculo contable institucional de las escuelas públicas.
        \CorrectChoice Son una herramienta en línea de Google que permite crear, editar y compartir textos o trabajos digitales desde cualquier dispositivo con acceso a internet.
        \choice Son aplicaciones locales de escritorio que guardan los archivos únicamente en el disco duro local de la computadora.
    \end{choices}
    

    \question ¿Qué herramientas básicas ofrece Documentos de Google para la correcta edición de un bloque de texto?
    \begin{choices}
        \choice Programación de macros automáticas y encriptación de seguridad militar para archivos de texto confidenciales.
        \choice Edición avanzada de audio digital y exportación directa a plataformas de streaming.
        \choice Herramientas exclusivas para el diseño arquitectónico y cálculo de planos bidimensionales.
        \CorrectChoice Cambiar tipo y tamaño de letra, aplicar negritas, cursivas y subrayado, insertar imágenes, tablas, enlaces, cambiar colores y alinear párrafos.
    \end{choices}
	
\newpage
\section{Ética, Naturaleza y Sociedades}


     \question En el modelo de transición demográfica, ¿qué fenómeno poblacional caracteriza a la segunda fase (transición incipiente) que experimentan los países en vías de desarrollo?
    \begin{choices}
        \choice Una alta tasa de natalidad acompañada de una alta tasa de mortalidad, manteniendo la población estable.
        \CorrectChoice Un descenso drástico en la tasa de mortalidad debido a mejoras médicas, mientras la natalidad se mantiene alta, generando un rápido crecimiento poblacional.
        \choice Una caída sostenida de la natalidad por debajo del nivel de reemplazo generacional, provocando envejecimiento poblacional.
        \choice Un equilibrio vegetativo donde la emigración supera a la inmigración, reduciendo la población absoluta.
    \end{choices}
	 

    \question El clima tipo "Af" (Ecuatorial lluvioso) según la clasificación de Köppen se asocia con un ecosistema específico que posee la mayor biodiversidad terrestre. ¿Qué factor físico determina esta condición?
    \begin{choices}
        \choice La marcada estacionalidad térmica que obliga a las especies a migrar anualmente.
        \choice La presencia de corrientes marinas frías que condensan la humedad en las costas.
        \CorrectChoice La constante convergencia de vientos alisios y la perpendicularidad de los rayos solares que generan altas temperaturas y precipitaciones todo el año.
        \choice La altitud extrema de los relieves montañosos que atrapan la humedad de los vientos monzónicos.
    \end{choices}
	

    \question ¿Cuál es la diferencia fundamental entre un riesgo natural y un desastre natural desde la perspectiva de la vulnerabilidad geográfica?
    \begin{choices}
        \choice El riesgo es predecible a corto plazo, mientras que el desastre es un fenómeno meteorológico aleatorio.
        \CorrectChoice El riesgo es la probabilidad de que un fenómeno extremo afecte un asentamiento, mientras que el desastre es la materialización del daño por la falta de preparación humana.
        \choice El riesgo aplica únicamente a fenómenos tectónicos (sismos), mientras que los desastres se refieren a eventos climáticos (huracanes).
        \choice No existe diferencia; ambos términos describen la magnitud de la energía liberada por la Tierra.
    \end{choices}
	

    \question En la conformación del espacio económico global, ¿qué función principal cumplen bloques económicos como el T-MEC o la Unión Europea?
    \begin{choices}
        \choice Estandarizar la moneda y el idioma de todos los países miembros para unificar la cultura.
        \choice Restringir totalmente la inversión extranjera proveniente de países de otros continentes.
        \choice Aislar económicamente a los países en desarrollo mediante bloqueos comerciales marítimos.
        \CorrectChoice Reducir o eliminar barreras arancelarias para facilitar el libre tránsito de mercancías, servicios y capitales entre las naciones que los integran.
    \end{choices}
	

    \question El concepto de "Huella Ecológica" es un indicador ambiental clave. ¿Qué mide exactamente este parámetro en una población determinada?
    \begin{choices}
        \choice La cantidad de emisiones de gases de efecto invernadero emitidas exclusivamente por el sector industrial.
        \CorrectChoice El área de territorio ecológicamente productivo necesaria para producir los recursos consumidos y asimilar los residuos generados por dicha población.
        \choice El porcentaje de especies endémicas que han sido desplazadas por la expansión de la mancha urbana.
        \choice La profundidad de los mantos acuíferos afectados por la contaminación de metales pesados.
    \end{choices}
	

    \subsection{Formación Cívica y Ética}

    \question Desde la perspectiva ética, ¿cuál es la diferencia principal entre la autonomía moral y la heteronomía?
    \begin{choices}
        \CorrectChoice En la autonomía el individuo actúa según principios racionales y convicciones propias; en la heteronomía, actúa sometido a normas impuestas por otros (sociedad, leyes, costumbres) sin cuestionarlas.
        \choice La autonomía es la capacidad de evadir las leyes constitucionales, mientras que la heteronomía es el cumplimiento estricto de las mismas.
        \choice En la autonomía se busca el beneficio de la comunidad, y en la heteronomía se busca exclusivamente el placer individual.
        \choice La autonomía es característica de los regímenes dictatoriales, y la heteronomía es la base de las repúblicas democráticas.
    \end{choices}
	

    \question La sexualidad humana es un concepto integral compuesto por cuatro potencialidades. ¿Cuáles son estas dimensiones?
    \begin{choices}
        \choice Anatómica, fisiológica, neurológica y genética.
        \CorrectChoice Reproductividad, género, erotismo y vínculos afectivos.
        \choice Biológica, legal, económica y demográfica.
        \choice Impulso, instinto, represión y sublimación.
    \end{choices}
	

    \question En México, ¿cuál es la función primordial de la Comisión Nacional de los Derechos Humanos (CNDH)?
    \begin{choices}
        \choice Dictar sentencias penales y encarcelar a los delincuentes de alta peligrosidad.
        \choice Modificar la Constitución Política para crear nuevos derechos laborales.
        \CorrectChoice Recibir quejas, investigar y emitir recomendaciones públicas no vinculantes sobre presuntas violaciones de derechos humanos cometidas por autoridades federales.
        \choice Representar a México en los juicios de la Corte Penal Internacional en La Haya.
    \end{choices}
	

    \question En la teoría de la resolución de conflictos, ¿qué diferencia existe entre la "paz negativa" y la "paz positiva"?
    \begin{choices}
        \choice La paz negativa se logra mediante tratados internacionales; la positiva mediante guerras de dominación.
        \CorrectChoice La paz negativa es simplemente la ausencia de guerra o violencia directa; la paz positiva implica la presencia de justicia social, equidad y respeto estructural a los derechos humanos.
        \choice La paz negativa es temporal y decretada por un juez; la paz positiva es perpetua y religiosa.
        \choice No existe distinción, ambos términos son sinónimos en el derecho internacional humanitario.
    \end{choices}
	

    \question El concepto de "Ciudadanía Digital" implica derechos y responsabilidades. ¿Qué práctica atenta directamente contra la ética de la convivencia en entornos virtuales?
    \begin{choices}
        \choice La participación en foros académicos de debate político abierto.
        \choice La protección de los datos personales mediante configuraciones de privacidad avanzadas.
        \CorrectChoice El \textit{cyberbullying} (ciberacoso), la difusión no consentida de información íntima y la proliferación de discursos de odio.
        \choice La validación y contraste de noticias (fact-checking) antes de compartirlas en redes sociales.
    \end{choices}
	

    \subsection{Historia}

    \question La Ilustración fue un movimiento intelectual del siglo XVIII. ¿Qué principio político, propuesto por Montesquieu, fue fundamental para el desarrollo de las democracias modernas y las revoluciones burguesas?
    \begin{choices}
        \choice La justificación del derecho divino de los reyes absolutistas.
        \choice El establecimiento de un sistema de castas basado en el nacimiento.
        \CorrectChoice La separación de los poderes del Estado en Ejecutivo, Legislativo y Judicial para evitar la tiranía.
        \choice La abolición total de la propiedad privada de los medios de producción.
    \end{choices}
	

    \question Durante la Segunda Revolución Industrial (fines del siglo XIX), ¿cuáles fueron las nuevas fuentes de energía que desplazaron parcialmente al carbón y la máquina de vapor?
    \begin{choices}
        \CorrectChoice La electricidad y los derivados del petróleo (motor de combustión interna).
        \choice La energía eólica y la energía solar termodinámica.
        \choice La energía nuclear de fisión y la biomasa.
        \choice La fuerza motriz animal y los molinos de marea.
    \end{choices}
	

    \question ¿Cuál de las siguientes opciones describe la principal causa estructural (no el detonante) que condujo al estallido de la Primera Guerra Mundial?
    \begin{choices}
        \choice El colapso repentino de la bolsa de valores de Wall Street y la Gran Depresión.
        \CorrectChoice El expansionismo imperialista y la intensa competencia geopolítica entre las potencias europeas por el control de colonias y mercados.
        \choice La invasión de Polonia por las tropas de la Alemania nazi bajo el mando de Adolf Hitler.
        \choice El surgimiento del comunismo en Rusia y la formación de la Unión Soviética.
    \end{choices}
	

    \question En el contexto de la Guerra Fría, ¿qué término define a los conflictos periféricos (como la Guerra de Corea o la Guerra de Vietnam) donde Estados Unidos y la URSS intervinieron indirectamente?
    \begin{choices}
        \choice Guerras santas o cruzadas contemporáneas.
        \choice Conflictos asimétricos de descolonización pacífica.
        \CorrectChoice Guerras subsidiarias o guerras "proxy" (por delegación).
        \choice Revoluciones burguesas de carácter industrial.
    \end{choices}
	

    \question La caída del Muro de Berlín en 1989 es un evento hito en la historia contemporánea porque simboliza:
    \begin{choices}
        \choice El inicio de la carrera espacial entre Estados Unidos y China.
        \CorrectChoice El fin de la división bipolar del mundo y el colapso inminente de los regímenes socialistas en Europa del Este.
        \choice El fracaso del Tratado de Versalles y el rearme alemán.
        \choice La consolidación militar del Pacto de Varsovia en el occidente europeo.
    \end{choices}
	
	
\end{questions}
\end{document}
